\section{Исследовательский раздел}

В данном разделе проводится исследование эффективности разработанного метода оценки кредитоспособности физических лиц на основе алгоритма градиентного бустинга деревьев решений (GBDT).

Цель исследования заключается в оценке качества разработанной модели и сравнении ее результатов с традиционными алгоритмами машинного обучения, применяемыми в задачах кредитного скоринга.

Для достижения поставленной цели решались следующие задачи:
\begin{itemize}
    \item анализ дисбаланса классов в исходном наборе данных;
    \item сравнение качества моделей логистической регрессии, случайного леса и градиентного бустинга;
    \item оценка качества моделей с использованием метрик ROC-AUC, Precision, Recall и F1-score;
    \item анализ влияния порога классификации на качество выявления дефолтов;
    \item анализ наиболее значимых признаков модели.
\end{itemize}

Home Credit Default Risk представляет собой открытый набор данных, размещенный на платформе Kaggle, и использовался в работе в качестве исходных данных. Набор данных содержит сведения о заявках клиентов, кредитной истории, предыдущих заявках и платежной дисциплине заемщиков. Положительным классом в рассматриваемой задаче является наступление дефолта заемщика, то есть значение целевой переменной TARGET = 1 \cite{KaggleHomeCreditDefaultRisk}.

Перед обучением моделей категориальные признаки были преобразованы методом one-hot encoding. Исходный набор данных был разделен на обучающую и тестовую выборки в соотношении 80\% и 20\%. Разделение выполнялось с сохранением соотношения классов с использованием стратификации. Оценка качества моделей проводилась на тестовой выборке, не использовавшейся при обучении.

Для учета дисбаланса классов при обучении модели LightGBM использовался параметр scale\_pos\_weight, рассчитываемый как отношение количества объектов отрицательного класса к количеству объектов положительного класса в обучающей выборке.

В качестве сравниваемых моделей были выбраны:
\begin{enumerate}
    \item Логистическая регрессия;
    \item Случайный лес (Random Forest);
    \item Градиентный бустинг (LightGBM).
\end{enumerate}

Выбранные модели позволяют сравнить линейный подход, ансамблевый метод на основе независимых деревьев и метод градиентного бустинга.

\subsection{Анализ дисбаланса классов}

Для исходного набора данных характерен существенный дисбаланс классов: доля клиентов с дефолтом значительно ниже доли надежных заемщиков. Доля положительного класса составляет около 8\%, что означает, что модель, всегда предсказывающая отсутствие дефолта, будет иметь высокое значение Accuracy, но не будет решать основную задачу выявления проблемных заемщиков.

В связи с этим метрика Accuracy не может рассматриваться как основная метрика качества, поскольку при выраженном дисбалансе классов она может давать завышенную оценку качества модели за счет преобладания объектов большинства класса. Для более корректной оценки моделей в таких условиях использовались метрики, учитывающие качество выявления положительного класса \cite{DavisGoadrichPRROC}:

\begin{itemize}
    \item ROC-AUC;
    \item Precision;
    \item Recall;
    \item F1-score;
    \item Precision-Recall кривая.
\end{itemize}

Особое внимание уделялось метрике Recall для положительного класса, поскольку в задаче кредитного скоринга пропуск заемщика с высоким риском дефолта может привести к финансовым потерям банка.

\subsection{Оптимизация гиперпараметров модели}

Для повышения качества разработанной модели была проведена оптимизация гиперпараметров алгоритма градиентного бустинга LightGBM. В рамках исследования рассматривались два подхода к подбору параметров: случайный поиск RandomizedSearchCV и полный перебор GridSearchCV.

RandomizedSearchCV выполняет проверку заданного количества случайно выбранных комбинаций гиперпараметров из общего пространства поиска. Такой подход позволяет сократить вычислительные затраты по сравнению с полным перебором, поскольку число проверяемых комбинаций задается заранее. Метод GridSearchCV, в свою очередь, выполняет исчерпывающий перебор всех комбинаций параметров из заданной сетки, что обеспечивает полный поиск по выбранному пространству, но требует значительно больше времени при увеличении числа параметров и их возможных значений.

В процессе оптимизации подбирались следующие параметры:
\begin{itemize}
    \item количество деревьев (n\_estimators);
    \item скорость обучения (learning\_rate);
    \item количество листьев (num\_leaves);
    \item максимальная глубина деревьев (max\_depth);
    \item минимальное количество объектов в листе (min\_child\_samples).
\end{itemize}

Подбор параметров осуществлялся с использованием трехкратной кросс-валидации с оценкой качества по метрике ROC-AUC, поскольку данная метрика оценивает способность модели ранжировать объекты разных классов и менее чувствительна к выбору конкретного порога классификации, чем Accuracy.

Для RandomizedSearchCV использовалось расширенное пространство гиперпараметров. Метод случайного поиска проверял ограниченное количество случайно выбранных комбинаций, что позволило выполнить подбор параметров за приемлемое время и избежать полного перебора всех возможных сочетаний.

В результате применения RandomizedSearchCV было получено значение F1-score, равное 0.2173 и ROC-AUC, равное 0.7758. Полученные результаты показывают, что случайный поиск позволил улучшить значение F1-score по сравнению с базовой конфигурацией модели.

Дополнительно был выполнен подбор гиперпараметров с использованием GridSearchCV на сокращенном пространстве параметров. Такой подход позволил выполнить полный перебор в рамках ограниченной сетки и сравнить его с результатами случайного поиска без чрезмерного увеличения времени вычислений.

В результате применения GridSearchCV на сокращенной сетке были получены следующие параметры:
\begin{itemize}
    \item learning\_rate = 0.03;
    \item max\_depth = 8;
    \item min\_child\_samples = 100;
    \item n\_estimators = 500;
    \item num\_leaves = 31.
\end{itemize}

Значение метрики ROC-AUC для GridSearchCV составило 0.7771, значение F1-score составило 0.2152. Таким образом, полный перебор на сокращенном пространстве параметров позволил получить наибольшее значение ROC-AUC среди рассмотренных вариантов, однако прирост относительно базовой модели оказался незначительным.

Для сравнения также была оценена базовая конфигурация модели без автоматического подбора гиперпараметров. Значение ROC-AUC для базовой модели составило 0.7757, значение F1-score составило 0.2098.

Дополнительно оценена вычислительная сложность полного перебора гиперпараметров на расширенном пространстве параметров. При использовании 108 комбинаций параметров и трехкратной кросс-валидации общее количество запусков обучения модели составило бы 324. С учетом экспериментально измеренного времени обучения одной комбинации параметров полный перебор потребовал бы значительно большего времени, чем случайный поиск.

GridSearchCV при расширенном пространстве параметров не является целесообразным решением в условиях ограниченных вычислительных ресурсов.
RandomizedSearchCV позволяет получить сопоставимое качество модели при меньшем количестве запусков обучения, что делает данный метод практически полезным при подборе гиперпараметров.

Результаты сравнения различных методов оптимизации гиперпараметров представлены на рисунке~\ref{fig:opt_comparison}.

\begin{figure}[ht]
    \centering
    \includegraphics[width=\linewidth]{inc/png/optimization_comparison.png}
    \caption{Сравнение методов оптимизации гиперпараметров}
    \label{fig:opt_comparison}
\end{figure}

По результатам, представленным на графике, значения основной метрики у всех рассмотренных методов оптимизации находятся на близком уровне. Базовая модель без оптимизации показывает значение 0.7757, GridSearchCV --- 0.7771, а RandomizedSearchCV --- 0.7758. Таким образом, полный перебор параметров позволил получить наибольшее значение ROC-AUC, однако улучшение относительно базовой модели оказалось незначительным.

Для метрики F1-score наилучший результат был получен при использовании RandomizedSearchCV и составил 0.2173. Базовая модель показала значение 0.2098, а GridSearchCV --- 0.2152. Это показывает, что случайный поиск позволил улучшить баланс между Precision и Recall по сравнению с базовой конфигурацией.

Таким образом, можно сделать вывод, что оба метода оптимизации гиперпараметров позволяют получить сопоставимое качество модели. GridSearchCV обеспечивает немного большее значение ROC-AUC, RandomizedSearchCV показывает лучшее значение F1-score и требует меньшего количества проверяемых комбинаций. Поэтому при ограниченных вычислительных ресурсах случайный поиск является более предпочтительным, тогда как полный перебор целесообразно использовать только на сокращенном пространстве параметров.


\subsection{Сравнение моделей по основным метрикам}

Для сравнения моделей использовались метрики ROC-AUC и F1-score. ROC-AUC показывает, насколько хорошо модель различает заемщиков с разным уровнем кредитного риска независимо от выбранного порога классификации. Метрика F1-score оценивает баланс между Precision и Recall для положительного класса, которым в данной задаче считается дефолт заемщика.

Результаты сравнения моделей по основным метрикам представлены на рисунке~\ref{fig:main_metrics}.

\begin{figure}[ht]
    \centering
    \includegraphics[width=\linewidth]{inc/png/main_metrics.png}
    \caption{Сравнение моделей по основным метрикам}
    \label{fig:main_metrics}
\end{figure}

По метрике ROC-AUC наилучший результат показала модель LightGBM. Значение ROC-AUC для нее составило 0.78. Логистическая регрессия показала значение 0.75, а случайный лес --- 0.73. Это означает, что модель LightGBM лучше остальных рассмотренных алгоритмов ранжирует заемщиков по вероятности дефолта.

По F1-score метрике лучший результат был получен для LightGBM. Значение F1-score составило 0.22. Для логистической регрессии значение F1-score составило 0.18, для случайного леса --- 0.14. Следовательно, при выбранном пороге классификации LightGBM обеспечивает более удачный баланс между выявлением дефолтных заемщиков и количеством ложноположительных срабатываний.

Таким образом, по совокупности ROC-AUC и F1-score модель LightGBM показала наиболее предпочтительный результат среди рассмотренных алгоритмов.

\subsection{Сравнение моделей по дополнительным метрикам}

Дополнительно были рассмотрены метрики Accuracy, Precision и Recall. Их значения представлены на рисунке~\ref{fig:extra_metrics}.

\begin{figure}[ht]
    \centering
    \includegraphics[width=\linewidth]{inc/png/extra_metrics.png}
    \caption{Сравнение моделей по дополнительным метрикам}
    \label{fig:extra_metrics}
\end{figure}

Высокие значения Accuracy были получены для логистической регрессии и случайного леса: 0.91 и 0.92 соответственно. Однако в данной задаче высокая Accuracy не является достаточным показателем качества модели. Это связано с дисбалансом классов: надежных заемщиков в выборке значительно больше, чем заемщиков с дефолтом. Поэтому модель может часто предсказывать отсутствие дефолта и получать высокий общий процент правильных ответов, но при этом плохо выявлять проблемных клиентов.

Значения Recall для логистической регрессии и случайного леса составили 0.12 и 0.09 соответственно. Это означает, что данные модели находят только небольшую часть заемщиков, фактически допустивших дефолт. Для задачи кредитного скоринга такой результат является недостаточным, поскольку пропуск рискованного заемщика может привести к финансовым потерям банка.

Модель LightGBM, напротив, показала существенно более высокое значение Recall, равное 0.89. Это означает, что модель выявляет большую часть заемщиков с высоким риском дефолта. При этом значение Accuracy у LightGBM составило 0.48, а Precision --- 0.12. Низкое значение Precision связано с тем, что при выбранном пороге классификации модель чаще относит заемщиков к группе риска, увеличивая количество ложноположительных решений.

Для задачи кредитного скоринга такой результат является ожидаемым: чем больше дефолтных заемщиков выявляет модель, тем чаще она ошибочно относит к группе риска надежных клиентов. В данной работе использовался порог классификации 0.3, что позволило повысить Recall и уменьшить количество пропущенных дефолтных заемщиков. Такой выбор порога оправдан тем, что для банка пропуск проблемного заемщика является более критичным, чем направление надежного клиента на дополнительную проверку.

Следовательно, метрика Accuracy не может рассматриваться как основная метрика качества в данной задаче. Более важными являются ROC-AUC, Recall и F1-score, так как они позволяют оценить способность модели работать с редким положительным классом.


\subsection{Анализ ROC-кривых}

Были построены ROC-кривые сравниваемых моделей для дополнительной оценки качества классификации. Данный график позволяет оценить поведение модели не при одном фиксированном пороге классификации, а при изменении порога в диапазоне от 0 до 1.

ROC-кривая показывает зависимость доли верно выявленных положительных объектов от доли ложноположительных срабатываний при изменении порога классификации. В рассматриваемой задаче положительным классом является дефолт заемщика. Диагональная линия на графике соответствует случайному классификатору, а расположение кривой выше данной диагонали указывает на более высокую способность модели различать классы.

ROC-кривые сравниваемых моделей представлены на рисунке~\ref{fig:roc}.

\begin{figure}[ht]
    \centering
    \includegraphics[width=\linewidth]{inc/png/roc_all_models.png}
    \caption{ROC-кривые моделей}
    \label{fig:roc}
\end{figure}

Из рисунка видно, что ROC-кривая модели LightGBM расположена выше кривых логистической регрессии и случайного леса на большей части диапазона значений False Positive Rate. Это соответствует наибольшему значению ROC-AUC, равному 0.777. Логистическая регрессия занимает промежуточное положение с ROC-AUC 0.755, а случайный лес показывает наименьшее значение ROC-AUC 0.729.

Полученные результаты подтверждают преимущество LightGBM. Данная модель лучше ранжирует заемщиков по уровню кредитного риска среди рассмотренных алгоритмов.

\subsection{Анализ Precision-Recall кривых}

Для детальной оценки качества моделей были построены Precision-Recall кривые. В отличие от ROC-кривой, данный график особенно полезен при несбалансированных классах, так как показывает качество модели именно при выявлении положительного класса.

В данной задаче положительный класс соответствует дефолту заемщика. Recall показывает долю выявленных дефолтных заемщиков, а Precision -- долю действительно дефолтных заемщиков среди тех, кого модель отнесла к рискованным.


Precision-Recall кривые представлены на рисунке~\ref{fig:pr}.

\begin{figure}[ht]
    \centering
    \includegraphics[width=\linewidth]{inc/png/pr_all_models.png}
    \caption{Precision-Recall кривые моделей}
    \label{fig:pr}
\end{figure}

Базовый уровень на графике соответствует доле положительного класса в выборке и составляет около 0.08. Расположение кривых выше данного уровня означает, что модели дают результат лучше случайного выбора заемщиков.

Модель LightGBM показывает наиболее устойчивый баланс между Recall и Precision. При увеличении Recall значение Precision у данной модели остается выше, чем у логистической регрессии и случайного леса на значительной части графика. Это означает, что LightGBM лучше подходит для выявления заемщиков с повышенным риском дефолта при разных значениях порога классификации.

Логистическая регрессия показывает результат, близкий к LightGBM на отдельных участках, однако в целом уступает ей по качеству ранжирования. Случайный лес располагается ниже остальных моделей на большей части диапазона Recall, что согласуется с его более низкими значениями ROC-AUC и F1-score.

\subsection{Анализ значимости признаков}

Для интерпретации результатов модели была проведена оценка значимости признаков LightGBM. Наиболее значимые признаки представлены на рисунке~\ref{fig:features}.

\begin{figure}[ht]
    \centering
    \includegraphics[width=\linewidth]{inc/png/feature_importance.png}
    \caption{Наиболее значимые признаки модели}
    \label{fig:features}
\end{figure}

На результат модели больше всего повлияли внешние скоринговые признаки EXT\_SOURCE\_1, EXT\_SOURCE\_2 и EXT\_SOURCE\_3. Высокая значимость данных признаков объясняется тем, что они уже содержат агрегированную информацию о надежности заемщика, полученную из внешних источников.

Значимыми факторами также являются возраст заемщика, сумма кредита, ежемесячный платеж, стаж занятости и признаки предыдущих заявок.
Эти признаки содержательно связаны с кредитным риском: возраст и занятость характеризуют социально-экономический профиль клиента, сумма кредита и платеж отражают долговую нагрузку, а история предыдущих заявок показывает прошлое кредитное поведение заемщика.

Полученные результаты подтверждают, что модель использует признаки, имеющие содержательную связь с предметной областью кредитного скоринга.

\subsection{Ограничения проведенного исследования}

В рамках проведенного исследования модель обучалась на статическом наборе данных, в котором признаки рассматриваются как фиксированные характеристики заемщика и его кредитной истории. Однако в реальных банковских системах кредитоспособность клиента может изменяться во времени.

К временным факторам относятся изменение дохода клиента, динамика кредитной нагрузки, появление новых заявок, изменение платежной дисциплины и возникновение просрочек. Учет таких изменений позволяет отслеживать ухудшение финансового состояния заемщика и повышать качество прогнозирования дефолта.

В рамках данной работы временные зависимости учитывались косвенно через агрегированные признаки, например показатели по предыдущим заявкам, кредитной истории и платежам. Однако полноценное моделирование временной динамики не проводилось.

Таким образом, одним из ограничений проведенного исследования является отсутствие полноценного учета временной динамики. Учет временных факторов может рассматриваться как направление дальнейшего развития разработанного метода.


\subsection*{Выводы}

В ходе исследования была проведена оценка качества разработанного метода оценки кредитоспособности физических лиц на основе алгоритма градиентного бустинга LightGBM. Для проверки эффективности метода были рассмотрены несколько моделей машинного обучения: логистическая регрессия, случайный лес и градиентный бустинг.

Результаты экспериментов показали, что модель LightGBM обеспечивает наилучшее качество среди рассмотренных алгоритмов по метрикам ROC-AUC и F1-score. Это свидетельствует о более высокой способности модели различать заемщиков с различным уровнем кредитного риска и более эффективно выявлять клиентов с вероятностью дефолта.

Было установлено, что в условиях дисбаланса классов метрика Accuracy не является достаточной для оценки качества модели. Логистическая регрессия и случайный лес показали высокие значения Accuracy, однако имели низкое значение Recall для положительного класса. Это означает, что данные модели пропускали значительную часть дефолтных заемщиков.

Использование пониженного порога классификации позволило повысить полноту выявления заемщиков с высоким риском дефолта. При этом снижение порога привело к уменьшению Precision, что связано с ростом числа ложноположительных решений. Такой результат отражает характерный для кредитного скоринга компромисс между снижением кредитного риска и сохранением числа одобренных заявок.

Анализ значимости признаков показал, что наибольшее влияние на результат модели оказывают внешние скоринговые оценки, возраст заемщика, сумма кредита, размер ежемесячного платежа, стаж занятости и признаки, связанные с предыдущими заявками. Данные признаки имеют содержательную связь с предметной областью и отражают основные факторы, влияющие на кредитный риск.

Таким образом, проведенное исследование подтвердило целесообразность использования алгоритма градиентного бустинга LightGBM для задачи оценки кредитоспособности физических лиц. Разработанный метод может применяться как инструмент предварительной оценки кредитного риска, при этом итоговое решение о выдаче кредита должно приниматься с учетом дополнительной проверки и экспертной оценки специалиста.